\documentclass[conference]{IEEEtran}
\IEEEoverridecommandlockouts
% The preceding line is only needed to identify funding in the first footnote. If that is unneeded, please comment it out.

% Türkçe karakter desteği
\usepackage[utf8]{inputenc}
\usepackage[T1]{fontenc}

\usepackage{cite}
\usepackage{amsmath,amssymb,amsfonts}
\usepackage{algorithmic}
\usepackage{graphicx}
\usepackage{textcomp}
\usepackage{xcolor}
\usepackage{subfiles} % Dosyaları bölmek için
\usepackage{todonotes}
\usepackage{booktabs}
\usepackage{url}
\usepackage{listings}
\usepackage{tikz}
\usetikzlibrary{shapes,arrows,positioning,calc}

\def\BibTeX{{\rm B\kern-.05em{\sc i\kern-.025em b}\kern-.08em
    T\kern-.1667em\lower.7ex\hbox{E}\kern-.125emX}}

% --- Renk Tanımlamaları ---
\definecolor{codegreen}{rgb}{0,0.6,0}
\definecolor{codegray}{rgb}{0.5,0.5,0.5}
\definecolor{codepurple}{rgb}{0.58,0,0.82}
\definecolor{backcolour}{rgb}{0.96,0.96,0.94}

% --- Listings Ayarları ---
\lstdefinestyle{mystyle}{
    backgroundcolor=\color{backcolour},   
    commentstyle=\color{codegreen},
    keywordstyle=\color{blue}\bfseries,
    numberstyle=\tiny\color{codegray},
    stringstyle=\color{codepurple},
    basicstyle=\ttfamily\scriptsize,
    breakatwhitespace=false,         
    breaklines=true,                 
    captionpos=b,                    
    keepspaces=true,                 
    numbers=left,                    
    numbersep=5pt,                  
    showspaces=false,                
    showstringspaces=false,
    showtabs=false,                  
    tabsize=4
}
\lstset{style=mystyle}

% --- Türkçe Karakter Düzeltmeleri ---
\addto\captionsturkish{%
  \renewcommand{\figurename}{Şekil}%
  \renewcommand{\tablename}{Tablo}%
  \renewcommand{\abstractname}{Özet}%
}

% --- Kısaltmalar ve Özel Makrolar (Acronyms & Custom Macros) ---
% Bu dosya main.tex'in önsözünde (preamble) % --- Kısaltmalar ve Özel Makrolar (Acronyms & Custom Macros) ---
% Bu dosya main.tex'in önsözünde (preamble) % --- Kısaltmalar ve Özel Makrolar (Acronyms & Custom Macros) ---
% Bu dosya main.tex'in önsözünde (preamble) \input{acronyms} ile çağrılmaktadır.
% LaTeX kuralları gereği, önsözde düz metin (text output) yazdırılması derleme hatasına yol açar.
% Bu nedenle buraya sadece komut tanımlamaları (\newcommand vb.) eklenmelidir.
% Örnek: \newcommand{\NOMA}{Non-Orthogonal Multiple Access}
 ile çağrılmaktadır.
% LaTeX kuralları gereği, önsözde düz metin (text output) yazdırılması derleme hatasına yol açar.
% Bu nedenle buraya sadece komut tanımlamaları (\newcommand vb.) eklenmelidir.
% Örnek: \newcommand{\NOMA}{Non-Orthogonal Multiple Access}
 ile çağrılmaktadır.
% LaTeX kuralları gereği, önsözde düz metin (text output) yazdırılması derleme hatasına yol açar.
% Bu nedenle buraya sadece komut tanımlamaları (\newcommand vb.) eklenmelidir.
% Örnek: \newcommand{\NOMA}{Non-Orthogonal Multiple Access}
 % acronyms.tex dosyasının projende var olduğundan emin ol

\begin{document}

% --- Makale Başlığı ---
% Şablonunuzdaki başlığı değiştirmek isterseniz aşağıdaki satırı yorumdan çıkarabilirsiniz:
% \title{USRP E310 \& LoRaWAN Sinyal Analiz Projesi}
\title{GNU Radio Tabanlı Güç Etki Alanlı NOMA Sistemi ile Ardışık Girişim İptali (SIC) Uygulaması}

% --- Yazar Bilgileri ---
\author{
    \IEEEauthorblockN{Eren Kale}
    \IEEEauthorblockA{\textit{Elektrik-Elektronik Mühendisliği} \\
    \textit{Marmara Üniversitesi}\\
    erenkale@marun.edu.tr}
    \and
    \IEEEauthorblockN{Armağan Bi}
    \IEEEauthorblockA{\textit{Elektrik-Elektronik Mühendisliği} \\
    \textit{Marmara Üniversitesi}\\
    armaganbi@marun.edu.tr}
}

\maketitle

% --- Özet (Abstract) ---
\begin{abstract}
Beşinci nesil (5G) ve sonrasındaki kablosuz haberleşme sistemlerinde spektral verimliliğin artırılması en kritik mühendislik hedeflerinden biri hâline gelmiştir. Bu bağlamda, Ortogonal Olmayan Çoklu Erişim (NOMA) tekniği, aynı zaman-frekans kaynağının birden fazla kullanıcı tarafından eş zamanlı paylaşılmasına olanak tanıyan bir fiziksel katman çözümü olarak öne çıkmaktadır. Bu çalışmada, Güç Etki Alanlı NOMA (PD-NOMA) mimarisi kullanılarak iki kullanıcılı bir inen hat (downlink) senaryosu için uçtan uca alıcı-verici sistemi tasarlanmış ve GNU Radio Companion (GRC) ortamında gerçeklenmiştir. Verici tarafında Süperpozisyon Kodlaması prensibi ile iki kullanıcının BPSK modüleli sinyalleri farklı güç katsayıları ($\alpha_1 = 0.894$, $\alpha_2 = 0.447$; sırasıyla \%80 ve \%20 güç tahsisi) ile birleştirilmiştir. Alıcı tarafında ise Ardışık Girişim İptali (SIC) algoritması uygulanarak zayıf kullanıcının sinyali güçlü kullanıcının sinyali çözülüp çıkarıldıktan sonra elde edilmiştir. Sistem, IEEE 802.11n standardına uygun LDPC kanal kodlaması ($n = 1296$, $k = 648$, $R = 1/2$) ve Diferansiyel BPSK (DBPSK) modülasyonu ile donatılmıştır. AWGN kanal modeli altında yapılan simülasyonlarda dosya transferinin doğruluğu kanıtlanmış olup, alıcıda kilitlenme (deadlock) ve sızıntı (bleeding) problemleri özel dairesel tamponlu Python blokları ile çözülmüştür.
\end{abstract}

% --- Anahtar Kelimeler (Keywords) ---
\begin{IEEEkeywords}
NOMA, Süperpozisyon Kodlaması, Ardışık Girişim İptali (SIC), GNU Radio, LDPC, DBPSK, SDR, Spektral Verimlilik.
\end{IEEEkeywords}

% --- Makale Bölümleri (Sections) ---
\section{Giriş ve Literatür Özeti}

\subsection{Motivasyon ve Problem Tanımı}
Kablosuz haberleşme sistemlerinin evrimi, birinci nesilden (1G) beşinci nesile (5G) ve ötesine uzanan süreçte, giderek artan veri trafiği, kullanıcı yoğunluğu ve düşük gecikme gereksinimlerini karşılama zorunluluğuyla şekillenmiştir. Uluslararası Telekomünikasyon Birliği (ITU) tarafından tanımlanan IMT-2020 vizyonuna göre, 5G sistemlerinin hedefleri arasında 10 Gbps tepe veri hızı, 1 ms gecikme süresi ve $10^6$ cihaz/$\text{km}^2$ bağlantı yoğunluğu yer almaktadır \cite{itu2015imt}. Bu hedefler doğrultusunda, sınırlı frekans spektrumunun verimli kullanılması kritik bir mühendislik problemi olarak öne çıkmaktadır.

Geleneksel kablosuz haberleşme sistemlerinde kullanılan Ortogonal Çoklu Erişim (Orthogonal Multiple Access - OMA) teknikleri (FDMA, TDMA, CDMA ve OFDMA), kullanıcıları frekans, zaman, kod veya alt-taşıyıcı boyutlarında birbirinden ayırarak ortogonallik sağlamaktadır. Bu ortogonallik ilkesi, kullanıcılar arası girişimi (inter-user interference) ortadan kaldırma avantajını sunarken, aynı zamanda kaynak tahsisinde esnekliği ve toplam spektral verimliliği sınırlandırmaktadır. Özellikle heterojen kullanıcı dağılımlarında, yani kanal koşullarının büyük farklılıklar gösterdiği senaryolarda, OMA yöntemlerinin Shannon kapasitesine yaklaşma kabiliyeti zayıflamaktadır.

\subsection{Non-Ortogonal Çoklu Erişim (NOMA) Kavramı}
NOMA, ortogonallik kısıtlamasını kaldırarak birden fazla kullanıcının \textbf{aynı zaman-frekans kaynağını} eş zamanlı olarak paylaşmasını mümkün kılan bir fiziksel katman çoklu erişim tekniğidir. Literatürde NOMA teknikleri iki ana kategoride sınıflandırılmaktadır:
\begin{enumerate}
    \item \textbf{Güç Etki Alanlı NOMA (Power-Domain NOMA - PD-NOMA):} Kullanıcılara farklı güç seviyeleri atanarak sinyallerin süperpozisyon kodlaması yoluyla üst üste bindirilmesine dayanır. Alıcı tarafında ardışık girişim iptali (Successive Interference Cancellation - SIC) algoritması ile girişim giderimi yapılır \cite{ding2017application}.
    \item \textbf{Kod Etki Alanlı NOMA (Code-Domain NOMA):} SCMA (Sparse Code Multiple Access) ve MUSA (Multi-User Shared Access) gibi teknikleri kapsayan, kullanıcıların seyrek veya düşük korelasyonlu kodlarla ayrıştırıldığı yaklaşımlardır.
\end{enumerate}
Bu projede, implementasyon karmaşıklığı ve kavramsal netlik açısından PD-NOMA mimarisi tercih edilmiştir.

\subsection{OMA ve NOMA Karşılaştırması}
NOMA'nın en belirgin üstünlüğü, aynı spektral kaynak üzerinde birden fazla kullanıcıya hizmet vermesi nedeniyle toplam sistem kapasitesini artırmasıdır. Yayın (broadcast) kanal kapasitesine ulaşmak için, güçlü kullanıcıya atanan düşük güç ve zayıf kullanıcıya atanan yüksek güç ile Shannon sınırına yaklaşılabilmektedir. Ancak, bu kazanım alıcı tarafındaki SIC karmaşıklığı pahasına elde edilmektedir.

Tablo~\ref{tab:oma_noma_comp} OMA ve NOMA tekniklerinin temel karakteristiklerini karşılaştırmaktadır.

\begin{table}[htbp]
\caption{OMA ve PD-NOMA Tekniklerinin Temel Karakteristikleri}
\label{tab:oma_noma_comp}
\centering
\begin{tabular}{p{2.3cm}p{2.6cm}p{2.7cm}}
\toprule
\textbf{Özellik} & \textbf{OMA (OFDMA vb.)} & \textbf{PD-NOMA} \\
\midrule
\textbf{Kaynak Tahsisi} & Her kullanıcıya ayrık alt-taşıyıcı/zaman dilimi & Tüm kullanıcılar aynı kaynağı paylaşır \\
\textbf{Kullanıcılar Arası Girişim} & Ortogonallik ile sıfırlanır & Kontrollü girişim, alıcıda SIC ile giderilir \\
\textbf{Spektral Verimlilik} & Kullanıcı sayısı ile ölçeklenmez & Kullanıcı sayısı arttıkça kapasite kazancı sunar \\
\textbf{Kullanıcı Adaleti} & Zayıf kullanıcılar aleyhine eşitsizlik & Güç tahsisi ile adalet kontrol edilebilir \\
\textbf{Alıcı Karmaşıklığı} & Düşük (doğrudan çözümleme) & Yüksek (SIC algoritması gereklidir) \\
\textbf{Kapasite Sınırı} & Sınırlı & Yayın (broadcast) kanal kapasitesine yaklaşır \\
\bottomrule
\end{tabular}
\end{table}

\subsection{Literatürde NOMA Çalışmaları}
NOMA kavramının teorik temelleri, çok kullanıcılı bilgi teorisi çerçevesinde Cover (1972) tarafından ortaya konulan yayın kanalı kapasitesi analizine dayanmaktadır \cite{cover1972broadcast}. Saito vd. (2013) tarafından yayımlanan çalışmada, güç etki alanlı NOMA'nın 3GPP-LTE sistemlerine entegrasyonu ilk kez önerilmiştir \cite{saito2013non}. Dai vd. (2015), NOMA'nın 5G kablosuz ağlar için bir anahtar teknoloji olduğunu ortaya koyan kapsamlı bir derleme çalışması sunmuştur \cite{dai2015non}. Ding vd. (2014), çeşitli kanal koşulları altında PD-NOMA'nın ergodik kapasitesini ve kesinti olasılığını (outage probability) analiz etmiştir \cite{ding2014performance}. İslam vd. (2017) ise güç tahsisi stratejileri ve SIC alıcı tasarımına ilişkin detaylı bir çerçeve geliştirmiştir \cite{islam2017power}.

LDPC kodlarının NOMA sistemlerine entegrasyonu açısından, Gallager (1962) tarafından ortaya konulan temel LDPC çerçevesi \cite{gallager1962low} ve MacKay ve Neal (1996) tarafından gösterilen Shannon sınırına yakın performans \cite{mackay1996near}, bu projede kullanılan IEEE 802.11n standardı LDPC kodlarının teorik altyapısını oluşturmaktadır \cite{ieee200980211n}.

GNU Radio platformu üzerinde NOMA implementasyonu açısından, SDR (Software Defined Radio) tabanlı çalışmalar sınırlı sayıdadır. Bu çalışma, uçtan uca bir PD-NOMA transceiver'ın GNU Radio ortamında IEEE 802.11n LDPC kodlama ve DBPSK modülasyonu ile birlikte gerçeklenmesi bakımından özgün bir katkı sunmaktadır.

\subsection{Projenin Amacı ve Kapsamı}
Bu projenin temel amacı, iki kullanıcılı inen hat PD-NOMA senaryosunda verici tarafında süperpozisyon kodlaması ile sinyal birleştirme, IEEE 802.11n LDPC kanal kodlaması ile hata dayanıklılığı sağlama, AWGN kanal modeli altında iletim simülasyonu ve alıcı tarafında çapraz korelasyon tabanlı SIC algoritması ile ardışık kullanıcı çözümleme adımlarını kapsayan uçtan uca bir haberleşme zincirini GNU Radio Companion ortamında tasarlamak, gerçeklemek ve performansını değerlendirmektir.

\section{Teorik Altyapı ve Matematiksel Model}

\subsection{Süperpozisyon Kodlaması}
Güç etki alanlı NOMA'nın temel mekanizması, birden fazla kullanıcının sinyallerinin farklı güç katsayılarıyla ağırlıklandırılarak toplam bir sinyal olarak iletilmesidir. İki kullanıcılı bir inen hat (downlink) senaryosunda, baz istasyonu (BS) tarafından iletilen süperpoze sinyal aşağıdaki şekilde ifade edilmektedir:
\begin{equation}
s(t) = \sqrt{a_1 P} \cdot x_1(t) + \sqrt{a_2 P} \cdot x_2(t)
\end{equation}
Burada, $P$ toplam iletim gücünü, $x_1(t)$ ve $x_2(t)$ sırasıyla Kullanıcı 1 (Uzak Kullanıcı) ve Kullanıcı 2 (Yakın Kullanıcı) için modüle edilmiş sembol dizilerini, $a_1$ ve $a_2$ ise güç tahsis katsayılarını ($a_1 + a_2 = 1$ ve $a_1 > a_2$) temsil etmektedir.

Uzak kullanıcının kanal sönümlemesi daha fazla olduğu için sinyal kalitesini korumak adına gücün büyük kısmı bu kullanıcıya tahsis edilir. Bu projede güç tahsisi $a_1 = 0.8$ ve $a_2 = 0.2$ olarak belirlenmiştir. Normalize edilmiş birim güç ($P = 1$) varsayımı altında, genlik katsayıları şu şekilde hesaplanır:
\begin{equation}
\alpha_1 = \sqrt{a_1} = \sqrt{0.8} \approx 0.894 \approx \frac{2}{\sqrt{5}}
\end{equation}
\begin{equation}
\alpha_2 = \sqrt{a_2} = \sqrt{0.2} \approx 0.447 \approx \frac{1}{\sqrt{5}}
\end{equation}
Bu katsayılar $\alpha_1^2 + \alpha_2^2 = 1.0$ güç korunumunu sağlamaktadır. Böylece süperpoze sinyal genlik cinsinden:
\begin{equation}
s(t) = 0.894 \cdot x_1(t) + 0.447 \cdot x_2(t)
\end{equation}
biçimini alır.

BPSK modülasyonu altında ($x_i \in \{-1, +1\}$), süperpoze sinyalin konstelasyon uzayı dört noktadan oluşmaktadır (Tablo~\ref{tab:constellation_analysis}).

\begin{table}[htbp]
\caption{BPSK Altında Süperpoze Konstelasyon Analizi}
\label{tab:constellation_analysis}
\centering
\begin{tabular}{cccc}
\toprule
$x_1$ & $x_2$ & $s = \alpha_1 x_1 + \alpha_2 x_2$ & Karar Bölgesi (User 1) \\
\midrule
+1 & +1 & $+0.894 + 0.447 = +1.341$ & Pozitif ($>0$) \\
+1 & -1 & $+0.894 - 0.447 = +0.447$ & Pozitif ($>0$) \\
-1 & +1 & $-0.894 + 0.447 = -0.447$ & Negatif ($<0$) \\
-1 & -1 & $-0.894 - 0.447 = -1.341$ & Negatif ($<0$) \\
\bottomrule
\end{tabular}
\end{table}

Görüleceği üzere $\alpha_1 > \alpha_2$ koşulu sağlandığı için, birleşik sinyal $s(t)$'nin işareti tamamen baskın olan $x_1(t)$ sembolü tarafından belirlenmektedir. Bu durum, alıcıda herhangi bir girişim giderimi yapmadan doğrudan sıfır eşik değeriyle $x_1(t)$ sinyalinin çözülebilmesini mümkün kılar.

\subsection{AWGN Kanal Modeli}
Verici çıkışından alıcı girişine kadar olan iletim kanalı, Toplamsal Beyaz Gauss Gürültüsü (AWGN) modeli ile temsil edilmektedir:
\begin{equation}
r(t) = h \cdot s(t) + n(t)
\end{equation}
Burada $h$ kanal katsayısını (bu çalışmada düz sönümleme için $h = 1.0$), $n(t)$ ise sıfır ortalamalı ve $\sigma_n^2$ varyanslı karmaşık Gauss gürültüsünü temsil eder. GNU Radio'daki \texttt{Channel Model} bloğunda konfigüre edilen gürültü gerilimi varsayılan olarak $\sigma_n = 0.1$ seçilmiştir. Bu değer $P=1$ iletim gücü altında $\text{SNR} = \frac{P}{\sigma_n^2} = 100 \approx 20\text{ dB}$ değerine karşılık gelir. Kanal modelinde ayrıca normalize frekans kayması ($\Delta f = 0.01$) ve zamanlama sapması ($\epsilon = 1.0001$) bozulmaları da simüle edilmektedir.

\subsection{Ardışık Girişim İptali (SIC)}
SIC, NOMA alıcısının temel bileşenidir ve çok kullanıcılı girişimin kademeli olarak giderilmesini sağlar. İki kullanıcılı senaryoda alıcıya ulaşan sinyal:
\begin{equation}
r(t) = \alpha_1 x_1(t) + \alpha_2 x_2(t) + n(t)
\end{equation}
şeklindedir. SIC süreci şu dört adımdan oluşur:
\begin{enumerate}
    \item \textbf{Güçlü Kullanıcının Çözülmesi:} Alınan sinyal $r(t)$ üzerinde, baskın güce sahip olan Kullanıcı 1'in sinyali doğrudan çözümlenir. Kullanıcı 2'nin sinyali bu aşamada girişim olarak ele alınır:
    \begin{equation}
    \hat{x}_1 = \text{Dec}\left\{ r(t) \right\} = \text{Dec}\left\{ \alpha_1 x_1(t) + \alpha_2 x_2(t) + n(t) \right\}
    \end{equation}
    Kullanıcı 1 için Sinyal-Girişim-Artı-Gürültü Oranı (SINR):
    \begin{equation}
    \text{SINR}_1 = \frac{a_1 P}{a_2 P + \sigma_n^2} = \frac{0.8}{0.2 + \sigma_n^2}
    \end{equation}
    \item \textbf{Güçlü Sinyalin Yeniden Sentezlenmesi:} Çözülen $\hat{x}_1$ verisi, verici tarafındaki işlem zinciriyle aynı adımlardan geçirilerek yeniden oluşturulur:
    \begin{equation}
    \hat{s}_1(t) = \alpha_1 \cdot \text{Mod}\left\{ \text{Enc}\left\{ \hat{x}_1 \right\} \right\}
    \end{equation}
    \item \textbf{Girişim İptali:} Yeniden oluşturulan güçlü kullanıcı sinyali, zamanlama ($\hat{\tau}$) ve faz ($\hat{\phi}$) kaymaları tahmin edilerek alınan ham sinyalden çıkarılır:
    \begin{equation}
    r_{\text{clean}}(t) = r(t) - \hat{s}_1(t - \hat{\tau}) \cdot \hat{A} \cdot e^{j\hat{\phi}}
    \end{equation}
    Burada $\hat{A}$ genlik düzeltme faktörüdür. Kod çözme hatasız ($\hat{x}_1 = x_1$) ve hizalama mükemmel ise girişim terimi tamamen sıfırlanır:
    \begin{equation}
    r_{\text{clean}}(t) = \alpha_2 x_2(t) + n(t)
    \end{equation}
    \item \textbf{Zayıf Kullanıcının Çözülmesi:} Temizlenmiş sinyal üzerinden yakın kullanıcının (User 2) sinyali çözümlenir:
    \begin{equation}
    \hat{x}_2 = \text{Dec}\left\{ r_{\text{clean}}(t) \right\}
    \end{equation}
    Kullanıcı 2 için elde edilen Sinyal-Gürültü Oranı (SNR) ideal durumda:
    \begin{equation}
    \text{SNR}_2 = \frac{a_2 P}{\sigma_n^2} = \frac{0.2}{\sigma_n^2}
    \end{equation}
\end{enumerate}

\subsubsection{SIC Hata Yayılımı}
SIC sürecinin başarısı, Adım 1'deki kod çözme doğruluğuna kritik biçimde bağlıdır. Kullanıcı 1'in hatalı çözülmesi durumunda artık girişim (residual interference) ortaya çıkar:
\begin{equation}
r_{\text{clean}}(t) = \alpha_2 x_2(t) + n(t) + \alpha_1 \left[ x_1(t) - \hat{x}_1(t) \right]
\end{equation}
$\alpha_1 \left[ x_1(t) - \hat{x}_1(t) \right]$ terimi hata yayılımını temsil etmekte olup, $\alpha_1 / \alpha_2 = 2$ olması nedeniyle her bir hatalı sembol, Kullanıcı 2'nin sinyalini $2:1$ oranında bastırma potansiyeline sahiptir.

Benzer şekilde, zamanlama hizalamasındaki 1 sembol kayma durumunda artık girişim:
\begin{equation}
r_{\text{clean}}[n] = \alpha_1(x_1[n] - x_1[n-1]) + \alpha_2 x_2[n] + n[n]
\end{equation}
biçimini alır. $\alpha_1 = 2\alpha_2$ olduğundan, bu diferansiyel girişim terimi Kullanıcı 2 sinyalini tamamen bastırabilir. Bu nedenle, güçlü kanal kodlaması (LDPC) ve hassas zamanlama hizalaması büyük önem taşımaktadır.

\subsection{LDPC Kanal Kodlaması}
Düşük Yoğunluklu Parite Kontrol (LDPC) kodları, Shannon sınırına yakın performans sunan modern ileri hata düzeltme (Forward Error Correction - FEC) kodlarıdır \cite{gallager1962low, mackay1996near}. Bu projede kullanılan LDPC kodu, IEEE 802.11n (Wi-Fi) standardından alınmış olup aşağıdaki parametrelere sahiptir:
\begin{itemize}
    \item Kod sözcüğü uzunluğu ($n$): $1296$ bit
    \item Mesaj uzunluğu ($k$): $648$ bit
    \item Parite bitleri ($n - k$): $648$ bit
    \item Kod oranı ($R = k/n$): $0.5$ (yarı oran)
    \item Parite kontrol matrisi: \texttt{n\_1296\_k\_0648\_ieee.alist}
    \item Alt-matris boyutu ($Z$): $54$ (yarı-döngüsel yapı)
\end{itemize}

Matrisin yarı-döngüsel (quasi-cyclic) yapısında parite kontrol matrisi $\mathbf{H}$, $Z = 54$ boyutlu döngüsel permütasyon alt-matrislerinden oluşmaktadır:
\begin{equation}
\mathbf{H} = \begin{bmatrix} \mathbf{P}_{0,0} & \mathbf{P}_{0,1} & \cdots & \mathbf{P}_{0,23} \\ \mathbf{P}_{1,0} & \mathbf{P}_{1,1} & \cdots & \mathbf{P}_{1,23} \\ \vdots & \vdots & \ddots & \vdots \\ \mathbf{P}_{11,0} & \mathbf{P}_{11,1} & \cdots & \mathbf{P}_{11,23} \end{bmatrix}
\end{equation}

Bu projede her paket $\text{payload\_size} = 77$ bayt ($= 616$ bit) kullanıcı verisine sahiptir. CRC-32 eklenmesi (4 bayt) sonrası $81$ bayt ($= 648$ bit) elde edilmekte olup, bu değer LDPC kodunun mesaj uzunluğu $k = 648$ ile birebir örtüşmektedir. LDPC kodlama sonrası her paket $n = 1296$ bit uzunluğundaki kod sözcüğüne dönüşür.

\subsection{Diferansiyel BPSK (DBPSK) Modülasyonu}
Diferansiyel kodlamada, bilgi verinin mutlak faz değerinde değil, ardışık semboller arasındaki faz farkında taşınır:
\begin{equation}
d_k = b_k \oplus d_{k-1}
\end{equation}
Burada $b_k$ giriş biti, $d_k$ diferansiyel kodlanmış bit ve $\oplus$ XOR işlemidir. BPSK sembol eşlemesi:
\begin{equation}
x_k = 2d_k - 1 = \begin{cases} +1 & d_k = 1 \\ -1 & d_k = 0 \end{cases}
\end{equation}
Diferansiyel kodlamanın avantajı, alıcıda mutlak faz referansı gerektirmeden çözümleme yapılabilmesidir; bu özellik, Costas döngüsünün $180^\circ$ faz belirsizliğini ortadan kaldırmaktadır.

Alıcı tarafında diferansiyel çözme, yumuşak karar (soft-decision) formatında gerçekleştirilmektedir:
\begin{equation}
\hat{b}_k^{\text{soft}} = -(y_k \cdot y_{k-1})
\end{equation}
Burada $y_k$ yumuşak konstelasyon çıkışıdır. Yumuşak karar bilgisinin korunması, aşağı yöndeki LDPC kod çözücünün performansı açısından kritik öneme sahiptir.

\subsection{Darbe Şekillendirme: Kök Yükseltilmiş Kosinüs (RRC) Filtresi}
Semboller arası girişimi (ISI) önlemek ve bant genişliğini kontrol etmek amacıyla, Kök Yükseltilmiş Kosinüs (RRC) darbe şekillendirme filtresi kullanılmıştır. RRC filtresinin frekans tepkisi:
\begin{equation}
H_{\text{RRC}}(f) = \begin{cases} \sqrt{T_s} & |f| \leq \frac{1-\beta}{2T_s} \\ \sqrt{\frac{T_s}{2}\left[1 + \cos\left(\frac{\pi T_s}{\beta}\left(|f| - \frac{1-\beta}{2T_s}\right)\right)\right]} & \frac{1-\beta}{2T_s} < |f| \leq \frac{1+\beta}{2T_s} \\ 0 & |f| > \frac{1+\beta}{2T_s} \end{cases}
\end{equation}
Burada $\beta$ aşırı bant genişliği (roll-off) faktörüdür. Bu projede $\beta = 0.35$ olarak konfigüre edilmiştir. Filtre uzunluğu $11 \times \text{sps} = 44$ dokunuş (tap) olarak belirlenmiştir.

\section{GNU Radio Tabanlı Sistem Tasarımı}
Bu bölümde, GNU Radio Companion (GRC) ortamında tasarlanan NOMA-SIC dosya transferi sisteminin tüm blokları, parametreleri ve sinyal akış yolları detaylı olarak incelenmektedir. Sistem, üç ana katmandan oluşmaktadır: Verici (Transmitter), Kanal Modeli (Channel Model) ve Alıcı (Receiver/SIC). Toplam blok envanteri 60'ın üzerindedir ve 16 sanal sinyal yolu (virtual sink/source) çifti ile karmaşık sinyal akışı yönetilmektedir.

\subsection{Sistem Değişkenleri}
Akış diyagramında tanımlanan global sistem değişkenleri Tablo~\ref{tab:system_variables}'de özetlenmiştir.

\begin{table}[htbp]
\caption{Global Sistem Değişkenleri}
\label{tab:system_variables}
\centering
\begin{tabular}{lll}
\toprule
\textbf{Değişken} & \textbf{Değer} & \textbf{Açıklama} \\
\midrule
\texttt{samp\_rate} & $200$ kHz & Örnekleme frekansı \\
\texttt{sps} & $4$ & Sembol başına örnek sayısı \\
\texttt{payload\_size} & $77$ bayt & Paket yük boyutu \\
\texttt{preamble\_size} & $250$ bayt & Preambül boyutu \\
\texttt{postamble\_size} & $8$ bayt & Postambül boyutu \\
\texttt{noise} & $0.1$ ($0$--$2$) & AWGN gürültü gerilimi \\
\texttt{freq\_offset} & $0.01$ ($\pm 0.25$) & Normalize frekans ofseti \\
\texttt{time\_offset} & $1.0001$ & Zamanlama sapma oranı \\
\texttt{constel} & BPSK & BPSK konstelasyon nesnesi \\
\texttt{preamble\_syms} & 64 sembol & BPSK preambül dizisi \\
\texttt{ldpc\_enc} & $R=1/2$ alist & LDPC kodlayıcı tanımı \\
\texttt{ldpc\_dec} & max\_iter: 50 & Kullanıcı 1 LDPC kod çözücü \\
\texttt{ldpc\_dec\_2} & max\_iter: 50 & Kullanıcı 2 LDPC kod çözücü \\
\bottomrule
\end{tabular}
\end{table}

Sembol hızı ve bant genişliği şu şekilde hesaplanır:
\begin{equation}
R_s = \frac{f_s}{\text{sps}} = \frac{200{,}000}{4} = 50{,}000 \text{ sembol/s}
\end{equation}
\begin{equation}
B = R_s \cdot (1 + \beta) = 50{,}000 \times 1.35 = 67{,}500 \text{ Hz}
\end{equation}

\subsection{Verici (Transmitter) Katmanı}
Verici katmanı, iki paralel kodlama zincirinden ve bir süperpozisyon birleştirme bloğundan oluşmaktadır.

\subsubsection{Kullanıcı 1 Verici Zinciri}
Kullanıcı 1'in verici zinciri aşağıdaki blok sırasını takip etmektedir:
\begin{itemize}
    \item \textbf{File Source:} \texttt{bpsk\_transmit.txt} dosyasındaki verileri bayt akışı olarak sisteme besler. \texttt{Repeat = True} ile sürekli iletim sağlanır.
    \item \textbf{Stream to Tagged Stream:} Akışı $77$ baytlık etiketlenmiş paketlere böler.
    \item \textbf{Stream CRC32:} Paketlere 4 baytlık CRC-32 ekleyerek boyutu $81$ bayta ($648$ bit $= k$) yükseltir.
    \item \textbf{Repack Bits (8$\to$1):} Bayt düzeyindeki verileri tek bit/bayt formatına dönüştürür.
    \item \textbf{Additive Scrambler:} LFSR tabanlı karıştırıcı (\texttt{Mask=0x8A}, \texttt{Seed=0x7F}, \texttt{Length=7}) veri dizisindeki uzun tekrarları kırarak DC sapmasını önler ve alıcı ile senkronizasyonu garanti eder.
    \item \textbf{FEC Extended Encoder:} 648 bitlik mesajı 1296 bitlik kod sözcüğüne ($n=1296$) dönüştürerek $R=0.5$ oranlı LDPC hata düzeltmesi sağlar.
    \item \textbf{Tagged Stream Multiply Length (x2):} Paket etiketini $1296$'ya günceller.
    \item \textbf{Protocol Formatter \& Tagged Stream Mux:} 250 baytlık preambül, başlık, LDPC yükü ve 8 baytlık postambülü birleştirerek çerçeve yapısını kurar.
    \item \textbf{Constellation Modulator (DBPSK):} Bitleri diferansiyel olarak kodlar, BPSK sembollerine eşler, $\text{sps}=4$ ve $\beta=0.35$ RRC filtresinden geçirir.
    \item \textbf{Tag Gate:} Akış etiketlerinin kanal modeline sızmasını önlemek için tüm etiketleri siler.
    \item \textbf{Multiply Const:} Genlik katsayısı $\alpha_1 = 0.894$ ile çarpılarak \%80 güç tahsisi yapılır.
\end{itemize}

\subsubsection{Kullanıcı 2 Verici Zinciri ve Birleştirme}
Kullanıcı 2'nin verici zinciri Kullanıcı 1 ile aynı yapıdadır, ancak dosya kaynağı olarak \texttt{bpsk\_transmit\_2.txt} kullanır ve genlik katsayısı $\alpha_2 = 0.447$ (\%20 güç tahsisi) ile çarpılır. İki kullanıcının sinyalleri \texttt{Add} bloğu ile toplanarak süperpoze NOMA sinyali $s(t) = 0.894 \cdot x_1(t) + 0.447 \cdot x_2(t)$ oluşturulur.

\subsection{Kanal Modeli (Channel Model)}
Süperpoze sinyal kanal modelinden geçirilir. Kanal modeli, gürültü gerilimi ($\sigma_n = 0.1$), normalize frekans kayması ($\Delta f = 0.01$) ve zamanlama sapması ($\epsilon = 1.0001$) bozulmalarını uygulayarak gerçekçi bir RF ortamı simüle eder.

\subsection{Alıcı (Receiver) Katmanı ve SIC Döngüsü}
Alıcı ön-ucunda eşleştirilmiş filtreleme (\texttt{FFT RRC Filter}), sembol zamanlama kurtarma (\texttt{Symbol Synchronizer} - \texttt{SIGNAL\_TIMES\_SLOPE\_ML} TED ve 8-taps MMSE interpolatörü) ve taşıyıcı faz kilitlenmesi (\texttt{Costas Loop}) gerçekleştirilir. \texttt{Correlation Estimator} preambül üzerinden paket sınırlarını tespit eder.

\subsubsection{Kullanıcı 1 Çözümleme ve Yeniden Kodlama}
\texttt{Constellation Soft Decoder} karmaşık sembolleri BPSK LLR yumuşak kararlarına dönüştürür. Geliştirilen \texttt{Soft Diff Decoder} (EPY Block 0) yumuşak kararları koruyarak diferansiyel kod çözümü yapar. Paket sınırları belirlendikten sonra \texttt{LDPC Decoder} 50 iterasyonla User 1 verisini çözer ve CRC-32 kontrolünden geçen paketler \texttt{bpsk\_receive.txt} dosyasına yazılır.

Çözülen Kullanıcı 1 bitleri, \texttt{LDPC Encoder} ile tekrar 1296 bite kodlanır, BPSK sembol uzayına ($\{-1, +1\}$) eşlenir ve $\alpha_1 = 0.894$ ile çarpılarak yeniden sentezlenir.

\subsubsection{SIC Çıkarma ve Kullanıcı 2 Çözümleme}
Yeniden sentezlenen User 1 sinyali, \texttt{Decoupled NOMA SIC Aligner} (EPY Block 1) vasıtasıyla alınan sinyalle çapraz korelasyon üzerinden hizalanarak composite sinyalden çıkarılır. Girişimi temizlenmiş sinyal, Kullanıcı 2 alıcı zincirine beslenerek çözümlenir ve \texttt{bpsk\_receive\_2.txt} dosyasına yazılır.

\subsection{Sanal Yönlendirme Haritası (Virtual Routing)}
Sinyal yollarını modüler hale getirmek için tanımlanan sanal yönlendirme noktaları Tablo~\ref{tab:virtual_routing} 'de özetlenmiştir.

\begin{table}[htbp]
\caption{Sanal Yönlendirme Portları ve Akış Haritası}
\label{tab:virtual_routing}
\centering
\begin{tabular}{lllp{2.8cm}}
\toprule
\textbf{Çıkış (Sink)} & \textbf{Giriş (Source)} & \textbf{Tip} & \textbf{Açıklama} \\
\midrule
\texttt{transmit} & \texttt{transmit} & Complex & Süperpozisyon sonrası verici çıkışı \\
\texttt{channel} & \texttt{channel} & Complex & Kanal modelinin gürültülü çıkış sinyali \\
\texttt{recovery} & \texttt{recovery} & Complex & Costas Loop çıkışı (User 1 çözücü girişi) \\
\texttt{recovery\_2} & \texttt{recovery\_2} & Complex & Correlation Estimator çıkışı (SIC hizalama ham girişi) \\
\texttt{user\_2} & \texttt{user\_2} & Byte & User 1 çözücü çıkışındaki temiz bit akışı \\
\texttt{user\_2\_sub} & \texttt{user\_2\_sub} & Complex & Yeniden sentezlenmiş User 1 sinyali \\
\bottomrule
\end{tabular}
\end{table}

\subsection{Çerçeve Yapısı (Frame Structure)}
Sistemdeki paket iletiminin kararlılığı, çerçevelerin zaman alanında düzenli yapılandırılmasıyla sağlanmaktadır. Paket çerçeve yapısı Tablo~\ref{tab:frame_structure} 'de sunulmuştur.

\begin{table}[htbp]
\caption{Paket Çerçeve Yapısı}
\label{tab:frame_structure}
\centering
\begin{tabular}{llll}
\toprule
\textbf{Segment} & \textbf{İçerik} & \textbf{Boyut} & \textbf{Sembol Sayısı} \\
\midrule
Preambül & \texttt{[0xC0, 0xAF]} tekrarlı & $250$ B & $256$ sembol \\
Başlık (Header) & Erişim kodu + yük uzunluğu & Değişken & $\approx 64$ sembol \\
Yük (Payload) & LDPC kodlu kullanıcı verisi & $162$ B & $1296$ sembol \\
Postambül & \texttt{[0xC0, 0xAF]} tekrarlı & $8$ B & $\approx 64$ sembol \\
\bottomrule
\end{tabular}
\end{table}

\section{Simülasyon ve Performans Değerlendirmesi}

\subsection{Simülasyon Ortamı ve Koşulları}
Sistem performansı, GNU Radio Companion 3.10.12+ platformu üzerinde yazılımsal geri-döngü (software loopback) simülasyonlarıyla değerlendirilmiştir. İletim kanalı, AWGN gürültüsü ($\sigma_n = 0.1$, SNR $\approx 20$ dB), normalize frekans kayması ($\Delta f = 0.01$) ve zamanlama sapması ($\epsilon = 1.0001$) bozulmalarını barındırmaktadır. Verici ve alıcı taraflarında metin dosyaları kullanılarak uçtan uca dosya transferi gerçekleştirilmiş ve verilerin doğruluğu bayt düzeyinde test edilmiştir.

\subsection{LDPC Kodlamanın SIC Hata Yayılımı Üzerindeki Rolü}
NOMA alıcısında SIC aşamasının başarısı, birinci aşamadaki (Kullanıcı 1) kod çözme doğruluğuna doğrudan bağlıdır. Birinci aşamada meydana gelecek tek bir bit hatası, çıkarma işleminde girişimi yok etmek yerine iki katına çıkararak sonraki aşamada hata yayılımına (error propagation) neden olur. Bu riski minimize etmek amacıyla sisteme entegre edilen IEEE 802.11n standardına uygun LDPC kodunun etkileri şunlardır:
\begin{itemize}
    \item \textbf{Hata Düzeltme Kapasitesi:} $R=0.5$ oranlı LDPC kodu, 50 iterasyonlu logaritmik olasılık oranları (LLR) tabanlı inanç yayılımı (belief propagation) algoritması sayesinde gürültülü kanallarda $\text{BER} < 10^{-6}$ seviyesine ulaşılmasını ve artık girişimin minimize edilmesini sağlar.
    \item \textbf{Yumuşak Karar Sinerjisi:} Alıcı zincirinde \texttt{Constellation Soft Decoder} $\to$ \texttt{Soft Diff Decoder} yolu ile korunan yumuşak karar (LLR) bilgisi, LDPC kod çözücünün tam potansiyeline ulaşmasını sağlar. Sert karar tabanlı bir alıcıya kıyasla, yumuşak karar LDPC kombinasyonu 2--3 dB SNR kazancı sunar.
    \item \textbf{CRC ile Çift Katmanlı Koruma:} CRC-32 kontrolü, kod çözme sonrası paket bütünlüğünü doğrulayarak hatalı paketlerin SIC yeniden kodlama zincirine girmesini engeller ve hata yayılımını sınırlandırır.
\end{itemize}

\subsection{Karşılaşılan Kritik Sorunlar ve Uygulanan Çözümler}

\subsubsection{Zamanlayıcı Kilitlenmesi (Scheduler Deadlock) ve \textit{basic\_block} Çözümü}
İlk aşamada SIC hizalama bloğu bir \texttt{sync\_block} olarak tasarlanmıştır. \texttt{sync\_block} yapısında, her iki giriş portundan da eşit miktarda örnek gelmesi zorunludur. Ancak:
\begin{itemize}
    \item Giriş 0 (\texttt{in\_rx}), kanal çıkışından gelen anlık veridir (gecikme $\approx 0$).
    \item Giriş 1 (\texttt{in\_tx1}), alıcının veriyi çözmesi, LDPC kodunu çözmesi, CRC doğrulaması yapması ve ardından bu veriyi yeniden kodlayıp modüle etmesiyle oluşur (en az 1 çerçeve $\approx 1680$ sembol gecikmesi).
\end{itemize}
Bu durumda \texttt{sync\_block}, Giriş 1'den veri beklerken Giriş 0'daki veriyi işlemez ve bekletir. Giriş 0 verisi işlenemediği için de LDPC çözücüye veri gitmez ve Giriş 1 hiçbir zaman veri üretemez. Bu dairesel kilitlenme nedeniyle sistem tamamen durmaktadır.

Bu dairesel bağımlılığı kırmak için SIC bloğu \texttt{gr.basic\_block} yapısına çevrilmiştir. Giriş verileri geldikleri anda bağımsız olarak consume edilmiş ve dairesel iç tamponlarda saklanmıştır. \texttt{forecast} fonksiyonu \texttt{[0, 0]} döndürecek şekilde aşırı yüklenerek kilitlenme sorunu \%100 çözülmüştür.

\subsubsection{Sinyal Sızıntısı (Signal Bleeding) ve 1-Sembol Kayma Düzeltmesi}
SIC çıkarma işleminde meydana gelecek 1 sembollük zaman kayması ($\tau = 1$), kalan sinyalde User 1 bileşeninin tamamen yok edilememesine neden olur:
\begin{equation}
r_{\text{clean}}[n] = \alpha_1 (x_1[n] - x_1[n-1]) + \alpha_2 x_2[n] + n[n]
\end{equation}
Güç katsayıları $\alpha_1 = 0.894$ ve $\alpha_2 = 0.447$ olduğundan, $\alpha_1 / \alpha_2 = 2$ oranına sahiptir. Bu nedenle, çıkarılamayan 1 sembollük User 1 kalıntı bileşeni, çözülmek istenen User 2 sinyalini genlik bazında tamamen bastırır ve karar eşiği sıfır olduğundan alıcı çıkışına User 1 verileri sızar.

Senkronizasyon etiketlerinin analizi sonucunda, \texttt{Correlation Estimator} bloğunun \texttt{mark\_delay=16} parametresinin çerçeve başlangıcında 16 sembollük bir yapay kayma ürettiği saptanmıştır. Optimum düzeltme katsayısı:
\begin{equation}
\text{T\_offset} + 64 + \text{shift} = \text{T\_offset} + 48 \implies \text{shift} = -16
\end{equation}
olarak hesaplanmıştır. SIC hizalama bloğundaki paket başlangıç hesabı, bu doğrultuda $-16$ sembol kaydırılarak kalibre edilmiş ve sinyal sızıntısı tamamen giderilmiştir.

\subsection{Geliştirilen Özel Python Blokları}
Sistemde kullanılan özel gömülü Python bloklarının kaynak kodları aşağıda sunulmuştur.

\subsubsection{Yumuşak Diferansiyel Çözücü (\texttt{NOMA\_epy\_block\_0.py})}
LLR yumuşak karar bilgisini koruyarak diferansiyel çözme yapan Python bloğunun kodu Liste~\ref{lst:soft_diff_decoder}'de verilmiştir.

\begin{lstlisting}[language=Python, label=lst:soft_diff_decoder, caption={Yumuşak Diferansiyel Çözücü Python Bloğu}]
import numpy as np
from gnuradio import gr

class blk(gr.sync_block):
    def __init__(self, modulus=2):
        if modulus not in (2, 4):
            raise ValueError("Modulus 2 (BPSK) veya 4 (QPSK) olmalidir.")
        self.modulus = int(modulus)
        gr.sync_block.__init__(
            self,
            name='Soft Diff Decoder',
            in_sig=[np.float32],
            out_sig=[np.float32]
        )
        if self.modulus == 4:
            self.set_output_multiple(2)
        if self.modulus == 2:
            self.prev = np.float32(1.0)
        else:
            self.prev_complex = complex(1.0, 0.0)

    def work(self, input_items, output_items):
        inp = input_items[0]
        out = output_items[0]
        n = len(inp)
        if n == 0:
            return 0
        if self.modulus == 2:
            extended = np.empty(n + 1, dtype=np.float32)
            extended[0] = self.prev
            extended[1:] = inp
            out[:] = -(extended[1:] * extended[:-1])
            self.prev = np.float32(inp[-1])
        else:
            num_symbols = n // 2
            I_vals = inp[0::2]
            Q_vals = inp[1::2]
            curr = (I_vals + 1j * Q_vals).astype(np.complex64)
            prev_arr = np.empty(num_symbols, dtype=np.complex64)
            prev_arr[0] = self.prev_complex
            prev_arr[1:] = curr[:-1]
            y = curr * np.conj(prev_arr)
            out[0::2] = -(y.real.astype(np.float32))
            out[1::2] = -(y.imag.astype(np.float32))
            self.prev_complex = complex(curr[-1])
        return n
\end{lstlisting}

\subsubsection{Bağımsız NOMA SIC Hizalama Bloğu (\texttt{NOMA\_epy\_block\_1.py})}
Dairesel tamponlama ve çapraz korelasyon tabanlı dinamik hizalamayı kilitlenmesiz şekilde uygulayan Python bloğunun kodu Liste~\ref{lst:noma_sic_aligner}'de verilmiştir.

\begin{lstlisting}[language=Python, label=lst:noma_sic_aligner, caption={Decoupled NOMA SIC Aligner Python Bloğu}]
import numpy as np
from gnuradio import gr
import scipy.signal

class blk(gr.basic_block):
    def __init__(self, sample_rate=200000.0, near_user_amplitude=0.864, search_window=128, payload_size=1296, payload_offset=64):
        gr.basic_block.__init__(
            self,
            name='Decoupled NOMA SIC Aligner',
            in_sig=[np.complex64, np.complex64],
            out_sig=[np.complex64]
        )
        self.sample_rate = sample_rate
        self.near_user_amplitude = near_user_amplitude
        self.search_window = search_window
        self.payload_size = payload_size
        self.payload_offset = payload_offset
        self.buffer_rx = np.array([], dtype=np.complex64)
        self.buffer_tx1 = np.array([], dtype=np.complex64)
        self.rx_processed_abs = 0
        self.subtracted_until_abs = 0
        self.next_packet_start_abs = 0
        self.pkt_counter = 0

    def forecast(self, noutput_items, ninputs):
        return [0] * ninputs

    def general_work(self, input_items, output_items):
        in_rx = input_items[0]
        in_tx1 = input_items[1]
        out_rx2 = output_items[0]
        if len(in_rx) > 0:
            self.buffer_rx = np.concatenate((self.buffer_rx, in_rx))
            self.consume(0, len(in_rx))
        if len(in_tx1) > 0:
            self.buffer_tx1 = np.concatenate((self.buffer_tx1, in_tx1))
            self.consume(1, len(in_tx1))
        MAX_RX_BUFFER = 50000
        if len(self.buffer_rx) > MAX_RX_BUFFER:
            excess = len(self.buffer_rx) - MAX_RX_BUFFER
            self.subtracted_until_abs = max(self.subtracted_until_abs, self.rx_processed_abs + excess)
        MAX_TX1_BUFFER = 5 * self.payload_size
        if len(self.buffer_tx1) > MAX_TX1_BUFFER:
            excess = len(self.buffer_tx1) - MAX_TX1_BUFFER
            self.buffer_tx1 = self.buffer_tx1[excess:]
        while len(self.buffer_tx1) >= self.payload_size:
            if self.pkt_counter == 0:
                SEARCH_START = 1000
                SEARCH_END = 3000
            else:
                expected_payload_start_abs = self.next_packet_start_abs + 2048
                start_rel = expected_payload_start_abs - self.rx_processed_abs
                SEARCH_START = start_rel - 128
                SEARCH_END = start_rel + 128
            SEARCH_START = max(0, SEARCH_START)
            if len(self.buffer_rx) < SEARCH_END + self.payload_size:
                break
            rx_search_chunk = self.buffer_rx[SEARCH_START : SEARCH_END + self.payload_size]
            tx1_chunk = self.buffer_tx1[:self.payload_size]
            tx1_chunk_norm = np.sign(tx1_chunk.real)
            tx1_chunk_diff = np.cumprod(-tx1_chunk_norm)
            corr = scipy.signal.correlate(rx_search_chunk, tx1_chunk_diff, mode='valid')
            if len(corr) > 0:
                best_idx = np.argmax(np.abs(corr))
                best_val = corr[best_idx]
                payload_start_rel = SEARCH_START + best_idx
                payload_start_abs = self.rx_processed_abs + payload_start_rel
                tx1_energy = np.sum(np.abs(tx1_chunk_diff)**2)
                if tx1_energy > 1e-6:
                    estimated_amp = np.abs(best_val) / tx1_energy
                else:
                    estimated_amp = self.near_user_amplitude
                phase_offset = np.angle(best_val)
                if estimated_amp >= 0.03:
                    if estimated_amp > 1.10:
                        amplitude_scale = estimated_amp / 1.5
                    else:
                        amplitude_scale = estimated_amp
                    interference_signal = tx1_chunk_diff * amplitude_scale * np.exp(1j * phase_offset)
                    self.buffer_rx[payload_start_rel : payload_start_rel + self.payload_size] -= interference_signal
                    self.pkt_counter += 1
                    self.next_packet_start_abs = payload_start_abs - 2048 + 3408
                    self.subtracted_until_abs = max(self.subtracted_until_abs, payload_start_abs - 2048 + 3408)
                    self.buffer_tx1 = self.buffer_tx1[self.payload_size:]
                    continue
            self.pkt_counter += 1
            self.next_packet_start_abs += 3408
            self.buffer_tx1 = self.buffer_tx1[self.payload_size:]
        safe_to_output = self.subtracted_until_abs - self.rx_processed_abs
        safe_to_output = min(safe_to_output, len(self.buffer_rx))
        safe_to_output = max(0, safe_to_output)
        n_out = len(out_rx2)
        limit = min(safe_to_output, n_out)
        if limit > 0:
            out_rx2[:limit] = self.buffer_rx[:limit]
            self.buffer_rx = self.buffer_rx[limit:]
            self.rx_processed_abs += limit
            return limit
        else:
            return 0
\end{lstlisting}

\section{Sonuç ve Gelecek Çalışmalar}

\subsection{Elde Edilen Bulguların Özeti}
Bu çalışmada, Güç Etki Alanlı NOMA (PD-NOMA) downlink transceiver mimarisinin GNU Radio Companion platformu üzerinde uçtan uca dosya transferi yapabilen kararlı bir prototipi başarıyla geliştirilmiştir. Projenin temel katkıları aşağıdaki şekilde özetlenebilir:
\begin{itemize}
    \item \textbf{Süperpozisyon Kodlaması ve İletim Zinciri:} İki kullanıcının DBPSK modüleli sinyalleri $\alpha_1 = 0.894$ (\%80 güç) ve $\alpha_2 = 0.447$ (\%20 güç) katsayıları ile başarıyla birleştirilmiştir. İletim zincirinde IEEE 802.11n standardına uygun LDPC kodlaması ($R=1/2$, $n=1296$), LFSR tabanlı karıştırıcı, CRC-32 ve $\beta=0.35$ RRC darbe şekillendirme filtreleri entegre edilmiştir.
    \item \textbf{Alıcı ve SIC Mekanizması:} Güçlü kullanıcının (User 1) yumuşak karar LLR değerleri korunarak kodunun çözülmesi, bu verinin yeniden kodlanıp sembollere eşlenerek composite sinyalden çıkarılması işlemleri gerçekleştirilmiştir. Geliştirilen \textit{basic\_block} tabanlı dairesel tamponlama yapısı ile GNU Radio zamanlayıcısının dairesel bağımlılık kaynaklı kilitlenmesi (scheduler deadlock) sorunu aşılmıştır.
    \item \textbf{Performans Doğrulaması:} AWGN kanalı ($\text{SNR} \approx 20$ dB, normalize frekans kayması $\Delta f = 0.01$, zamanlama sapması $\epsilon = 1.0001$) altında her iki kullanıcı için başarılı metin dosyası transferi gerçekleştirilmiş ve verilerin bütünlüğü bayt düzeyinde otomatik olarak doğrulanmıştır.
\end{itemize}

\subsection{Projenin Akademik ve Mühendislik Katkısı}
Bu çalışma, NOMA'nın teorik sınırlarının ötesine geçerek, yazılım tabanlı radyo (SDR) platformlarında çalışabilen uçtan uca pratik bir demonstrasyon sunmaktadır. Özellikle dairesel akış geri beslemelerine sahip alıcı yapılarında zamanlayıcı kilitlenmesini çözen dairesel tamponlama mantığı, çapraz korelasyon tabanlı dinamik genlik/faz kalibrasyonu ve 1-sembollük sızıntı sapmasının giderilmesi projenin getirdiği özgün mühendislik çözümleridir.

\subsection{Gelecek Çalışmalar}
Proje sonuçları, gelecekte aşağıdaki yönlerde genişletilebilir:
\begin{enumerate}
    \item \textbf{USRP E310 Donanım Entegrasyonu:} Tasarlanan simülasyon modelinin, projedeki \texttt{Bpsk\_file\_transfer\_loopback.grc} akış diyagramı temel alınarak USRP E310 SDR donanımları üzerinden gerçek zamanlı 2.435 GHz bandında test edilmesi ve gerçek kanal sönümlemelerinin SIC performansına etkisinin incelenmesi.
    \item \textbf{PlutoSDR Entegrasyonu:} Analog Devices ADALM-PlutoSDR gibi düşük maliyetli SDR platformları kullanılarak laboratuvar ortamı kablosuz testlerinin yapılması.
    \item \textbf{Çok Kullanıcılı Genişletme:} İki kullanıcılı yapının üç veya daha fazla kullanıcıya genişletilerek kaskat SIC (kademeli girişim çıkarma) işleminin hata yayılımı üzerindeki etkilerinin analiz edilmesi.
    \item \textbf{Uyarlanır Güç Tahsisi:} Kullanıcıların anlık kanal durum bilgilerine (CSI) göre $\alpha_1$ ve $\alpha_2$ güç katsayılarının dinamik olarak optimize edilmesi.
    \item \textbf{OFDM-NOMA Entegrasyonu:} BPSK tabanlı tek taşıyıcılı sistemden çok taşıyıcılı OFDM-NOMA mimarisine geçiş yapılarak frekans-seçici kanallarda çalışabilme kabiliyetinin kazandırılması.
\end{enumerate}


% --- Kaynakça ---
\bibliographystyle{IEEEtran}
\bibliography{references}

\end{document}
